\chapter{Felhasználói dokumentáció}
\label{ch:user}

Ez a fejezet az alkalmazás teljes leírását nyújtja a felhasználó számára. A probléma összefoglalásával kezdődik, ismerteti az alkalmazott módszereket, majd útmutatást ad a program futtatásához és a vele való interakcióhoz.

\section{A probléma meghatározása}

A Rush Hour rejtvény~\cite{rushhour_puzzle} egy tologatós játék, amelyet egy $6 \times 6$-os négyzetrácsos táblán játszanak. A tábla több féle járművet tartalmazhat: $1 \times 2$ mezőt elfoglaló \textbf{autókat} és $1 \times 3$ mezőt elfoglaló \textbf{buszokat}. Minden jármű csak előre és hátra, a hosszabbik tengelye mentén mozoghat, és nem tudják egymást elmozdítani. A cél a kijelölt (piros színű) \textbf{főautó} vízszintes elcsúsztatása, amíg az el nem éri a tábla jobb szélén található kijáratot (3. sor, a tábla jobb oldalán).

A rejtvény akkor tekinthető \textbf{megoldottnak}, ha a főautó orra közvetlenül a kijáratnál helyezkedik el. A játék állapottere az összes a kezdőállapotból elérhető konfigurációból áll, az állapotok közötti átmenetek pedig egy-egy jármű egy mezővel való elmozdításának felelnek meg a megengedett irányban.

Az alkalmazás két egymáshoz kapcsolódó kihívásra fókuszál:
\begin{enumerate}
    \item Keresőalgoritmusok implementálása és interaktív vezérlése a Rush Hour rejtvényen, megengedve a felhasználónak hogy tetszőlegesen megállítsa, léptesse, folytassa az algoritmus futtatását, és közben minden döntést ki tudjon elemezni.
    \item Az állapotgráf vizualizációja, ahol a csúcsok a felfedezett állapotok, az élek pedig a megengedett lépések egy háromdimenziós színtérben, erőalapú elrendezéssel (force-directed layout) szimulálva, hogy a szerkezetileg összefüggő állapotok közel kerüljenek egymáshoz.
\end{enumerate}


\section{A módszerek áttekintése}

Az alkalmazás három keresőalgoritmust valósít meg, amelyek mindegyike heurisztikus vezérlési stratégiát alkalmaz konfigurálható heurisztikus függvényekkel:

\begin{description}
    \item[Tabu keresés] Egy nemmódosítható lokális keresési módszer, amely egy heurisztika alapján választja ki a legjobb szomszédos állapotot, miközben memóriát tart fenn a legutóbb meglátogatott állapotokról (\emph{tabu lista}), hogy elkerülje a köröket.
    \item[Visszalépéses (Backtracking) keresés] Módosítható stratégiájú keresés, amely heurisztika által vezérelt sorrendben tárja fel az állapotokat a köröket elkerülve, zsákutca esetén visszalép, véges állapotgráfra garantálja a leállást.
    \item[A\textsuperscript{C} gráfkeresés] Módosítható stratégiájú gráfkeresés, amely az $f(n) = g(n) + h(n)$ kiértékelő függvényt használja, ahol $g(n)$ a mélység (a kezdőállapottól megtett lépések száma), $h(n)$ pedig a heurisztikus érték. Megengedhető és monoton heurisztikákkal ez az algoritmus megtalálja az optimális megoldást.
\end{description}

Négy heurisztikus függvény áll rendelkezésre, amelyek mindegyike nemnegatív, megengedhető és monoton~\cite{gregorics_grafker_heuristic}:
\begin{compactitem}
    \item \textbf{Null heurisztika} ($h = 0$): nem ad útmutatást; a \enquote{brute-force} viselkedés megfigyelésére hasznos.
    \item \textbf{Távolság heurisztika}: a kijuttatandó autó távolsága a kijárattól (tartománya: 0--4).
    \item \textbf{Távolság + blokkolás heurisztika}: a távolság, plusz egy költség minden egyes, a főautó és a kijárat között lévő foglalt mezőért.
    \item \textbf{Legkevesebb lépés heurisztika}: Még erősebb alsó becslést ad a megoldásig szükséges lépésekre: a távolság, plusz a főautó útjába álló járművek elmozdításához szükséges minimális lépésszámot.
\end{compactitem}

A megengedhetőség azt jelenti, hogy a heurisztika sosem becsüli túl a tényleges minimális költséget, a monotonitás pedig azt, hogy egy lépés legfeljebb 1-gyel változtathatja meg a heurisztika értékét. Ezek a tulajdonságok biztosítják, hogy az A$^C$ gráfkeresés optimális megoldást találjon~\cite{gregorics_grafker_heuristic}.

A 3D-s állapotgráf elrendezéséhez erőalapú megközelítést alkalmazunk: az élek rugókként működnek, amelyek adott távolság felé húzzák a szomszédos állapotokat, míg minden csúcspár egymásra a távolságuk négyzetével fordítottan arányos taszítóerőt fejtenek ki egymásra. A számításigényes $O(V^2)$ komplexitású taszítást a Barnes–Hut algoritmus~\cite{barneshut1986} nyolcasfa (octree) térbeli adatszerkezet implementálásával approximáljuk, és ezzel az idő komplexitását le tudjuk csökkenteni $O(V \log V)$-re.


\section{Rendszerkövetelmények}

Az alkalmazás önállóan futtatható fájlként érhető el Windowson és Linuxon, és nem igényel további szoftvertelepítést.
Van lehetőség a projektet a Godot editorből futtatni és exportálni futtatható fájl formátumba egyéb platformokra.
Ajánlott minimum képernyőfelbontás: $1280 \times 720$

\begin{table}[H]
    \centering
    \begin{tabular}{ | m{0.25\textwidth} | m{0.65\textwidth} | }
        \hline
        \textbf{Platform} & \textbf{Futtatás} \\
        \hline \hline
        \emph{Windows} & Futtassa közvetlenül a mellékelt \texttt{rushhour.exe} fájlt. \\
        \hline
        \emph{Linux} & Futtassa közvetlenül a mellékelt \texttt{rushhour.x86\_64} fájlt. \\
        \hline
        \emph{MacOS / egyéb} & Töltse le és csomagolja ki a \texttt{rushhour.zip} fájlt és nyissa meg (Import project) a kicsomagolt \texttt{project.godot} fájlt a Godot 4.6 Mono szerkesztővel, majd jobb felüli gombbal futtassa. \\
        \hline
    \end{tabular}
    \caption{Támogatott platformok és futtatási utasításaik}
    \label{tab:platforms}
\end{table}

A projekt szerkesztőben való megnyitásához szükséges előfeltételek:
\begin{compactitem}
    \item Godot 4.6 Mono (C\#) szerkesztő letöltésre került.~\cite{godot_engine}
    \item .NET SDK 9.0 vagy újabb verzió telepítve van.~\cite{dotnet9}
\end{compactitem}

\begin{figure}[H]
    \centering
    \begin{subfigure}[t]{0.35\textwidth}
        \centering
        \includegraphics[width=\linewidth]{import_project}
        \caption{Projekt importálása}
    \end{subfigure}\hfill
    \begin{subfigure}[t]{0.4\textwidth}
        \centering
        \includegraphics[width=\linewidth]{open_project}
        \caption{Projekt megnyitása}
    \end{subfigure}\hfill
    \begin{subfigure}[t]{0.2\textwidth}
        \centering
        \includegraphics[width=\linewidth]{run_project}
        \caption{Projekt futtatása}
    \end{subfigure}
    \caption{A projekt megnyitása és futtatása a Godot szerkesztőből}
    \label{fig:open-run-project}
\end{figure}


A forrásból való fordításhoz, ha az előfeltételek adottak, a projekt megnyitása után bal felül Project > Export opciónál található, hozzá kell adni egy exportálási konfigurációt a kiválasztott platformra, majd le kell tölteni még egy exportálási könyvtárat amit fel fog ajánlani a szerkesztő.

\begin{figure}[H]
    \centering
    \begin{subfigure}[t]{0.32\textwidth}
        \centering
        \includegraphics[width=\linewidth]{export_project}
        \caption{Exportálás megnyitása}
    \end{subfigure}\hfill
    \begin{subfigure}[t]{0.32\textwidth}
        \centering
        \includegraphics[width=\linewidth]{add_export_option}
        \caption{Export konfiguráció hozzáadása}
    \end{subfigure}\hfill
    \begin{subfigure}[t]{0.32\textwidth}
        \centering
        \includegraphics[width=\linewidth]{manage_export_templates}
        \caption{Exportálási sablonok kezelése}
    \end{subfigure}

    \vspace{0.5em}

    \begin{subfigure}[t]{0.32\textwidth}
        \centering
        \includegraphics[width=\linewidth]{go_online}
        \caption{Online módba lépés}
    \end{subfigure}\hfill
    \begin{subfigure}[t]{0.32\textwidth}
        \centering
        \includegraphics[width=\linewidth]{download_export_templates}
        \caption{Export sablonok letöltése}
    \end{subfigure}\hfill
    \begin{subfigure}[t]{0.32\textwidth}
        \centering
        \includegraphics[width=\linewidth]{export_project_last_step}
        \caption{Projekt exportálása}
    \end{subfigure}
    \caption{A projekt exportálásának lépései a Godot szerkesztőből}
    \label{fig:export-project}
\end{figure}

\section{Az alkalmazás áttekintése}

Az alkalmazás ablaka két fő területre oszlik:
\begin{compactitem}
    \item \textbf{Bal oldali panel}: a 3D-s állapotgráf vizualizációs színtere. 
    \item \textbf{Jobb oldali panel}: A konfigurációs és vezérlési opciók 3 fülre osztva.
\end{compactitem}


\begin{figure}[H]
    \centering
    \includegraphics[width=0.75\textwidth]{app_preview}
    \caption{Az alkalmazás egy algoritmus futása közben}
    \label{fig:app-preview}
\end{figure}

A bal oldali panelen belül található:

\begin{compactitem}
    \item Bal alul az állapotgráffal való interakciónak az irányításai
    \item Jobb felül a jelenleg futó algoritmushoz kapcsolódó információk:
    \begin{compactitem}
        \item Az algoritmus jelenlegi státusza
        \item Az algoritmus által megtett lépések száma
        \item A megtalált megoldás lépésszáma
        \item A jelenleg táblán mutatott állás heurisztikus értéke
    \end{compactitem}
\end{compactitem}

A Jobb oldali panel fülekre osztva:
\begin{compactitem}
		\item \textbf{Feladvány választó fül}: Kattintással lehet választani az előnézettel felsorolt feladványok közül. 
		\item \textbf{Keresőalgoritmus beállítások fül}: Indítható az állapotok felfedezése, beállítható, hogy milyen keresőalgoritmust, milyen heurisztikát szeretnénk alkalmazni a kiválasztott feladványra.
		\item \textbf{Játéktábla fül}: Tartalmazza a játék tábláját és a manuális és algoritmus mód választó gombokat, az algoritmust vezérlő: megállítás, folytatás, léptetés, újraindítás gombokat, egy csúszkát az algoritmus léptetési sebességéhez és kapcsolókat az állapotgráf megjelenéséhez kapcsolódóan.
	\end{compactitem}
\section{Szint betöltése}

Az alkalmazás különböző nehézségű, kiválogatott Rush Hour rejtvények gyűjteményét tartalmazza. Az összes rejtvény görgetve tekinthető meg.

Bármikor másik feladvány választható az alábbi módon:
\begin{enumerate}
    \item Navigáljon a jobb felül a füleken belül a \textbf{Levels} fülre.
    \item Kattintson a kívánt rejtvénynek megfelelő gombra.
    \item A játéktábla betölti a kezdeti konfigurációt, a 3D-s állapotgráf térben pedig megjelenik az a csúcs, amely a kezdőállapotot reprezentálja.
\end{enumerate}

\section{Kézi játék}

A rejtvény manuálisan is játszható a játéktáblával való közvetlen interakció útján:

\begin{enumerate}
    \item Minden mozgatható jármű mellett piros nyíl indikátorok jelennek meg, mutatva a lehetséges mozgási irányokat.
    \item Kattintson egy nyílra a megfelelő jármű egy mezővel való eltolásához.
    \item Minden kézi lépés egy új állapotot, valamint egy hozzá tartozó csúcsot és élt hoz létre a 3D-s gráfon, lehetővé téve a felhasználó számára az állapotgráf interaktív felépítését.
\end{enumerate}


A manuális módba bármikor át lehet kapcsolni a \textbf{Manual} gombra nyomással, vagy egy már felfedezett állapot csúcsára kattintással. Mindkét esetben a tábla a kiválasztott állapotot jeleníti meg és manuális módba vált. Ha ekkor épp fut keresőalgoritmus, az a háttérben tovább fut.
\begin{figure}[H]
    \centering
    \begin{subfigure}[t]{0.32\textwidth}
        \centering
        \includegraphics[width=\linewidth]{manual_moves}
        \caption{A feladvány megoldása manuális lépésekkel keresőalgoritmus alkalmazása nélkül}
    \end{subfigure}\hfill
    \begin{subfigure}[t]{0.32\textwidth}
        \centering
        \includegraphics[width=\linewidth]{manual_button}
        \caption{A tábla már nem követi az algoritmus lépéseit manuális módba váltás után.}
    \end{subfigure}\hfill
    \begin{subfigure}[t]{0.32\textwidth}
        \centering
        \includegraphics[width=\linewidth]{clicked_state}
        \caption{Állapot kiválasztása egy felfedezett csúcsra kattintás esetén}
    \end{subfigure}
    \caption{Manuális játék és csúcsra kattintás}
    \label{fig:manual-play}
\end{figure}

\section{A keresőalgoritmus konfigurálása}

A keresőalgoritmus futtatása előtt a felhasználó a \textbf{Solver settings} fülön konfigurálhatja a paramétereket:

\begin{description}
    \item[Algoritmus választása] Válasszon a Tabu keresés, Visszalépéses keresés vagy A\textsuperscript{C} gráfkeresés közül.
    \item[Tabu méret] (Csak Tabu keresésnél) Egy $n \geq 0$ egész szám, amely meghatározza, hány legutóbb meglátogatott állapotra emlékezzen a rendszer.
    \begin{compactitem}
        \item $n = 0$: mohó algoritmus --- memória nélkül mindig a legjobb szomszédra lép.
        \item $n = 1$: hegymászó algoritmus --- csak az előző állapotra emlékszik, és elkerüli az oda való visszalépést.
        \item $n > 1$: elkerüli az utolsó $n$ meglátogatott állapotot. Ha minden lehetséges lépés tabu állapotba vezetne, az algoritmus megoldás nélkül terminál.
    \end{compactitem}
    \item[Heurisztikus függvény] Válasszon a négy elérhető heurisztika közül (Null, Distance, Distance + blocks, Least-moves).
    \item[Véletlenszerűség kapcsoló] Ha engedélyezve van, az azonos heurisztikus pontszámú állapotok közül az algoritmus véletlenszerűen választ a determinisztikus választás helyett, ami ismételt futtatásoknál különböző bejárási utakat eredményezhet.
    \item[Alkalmaz (Apply) gomb] Nyomja meg, hogy a beállítások szerinti keresőalgoritmus inicializálva legyen a betöltött feladványra.
\end{description}

\begin{figure}[H]
    \centering
    \includegraphics[width=0.5\textwidth]{solver_settings}
    \caption{Keresőalgoritmusok beállítása és alkalmazása}
    \label{fig:solver-settings}
\end{figure}


\section{Keresőalgoritmusok futtatása}

A szint betöltése és a megoldó konfigurálás és alkalmazás (Apply gomb) után az algoritmus az \textbf{Game board} fülről vezérelhető:

\begin{description}
    \item[Lejátszás/Szünet] Az algoritmus folyamatos futtatásának indítása vagy megállítása.
    \item[Léptetés (Step)] Az algoritmus egyetlen lépésének végrehajtása, ami hasznos a részletes megfigyeléshez.
    \item[Újraindítás] Az algoritmus visszaállítása a kezdőállapotba a jelenlegi beállítások megtartásával.
    \item[Sebességcsúszka] Az algoritmus lépései közötti késleltetés beállítása. A csúszka jobbra tolásával a sebesség nő; a maximum állásban az algoritmus a lehető leggyorsabban fut, képkockánként több száz lépést is végrehajthat.
\end{description}

\begin{figure}[H]
    \centering
    \includegraphics[width=0.5\textwidth]{algorithm_controls}
    \caption{Algoritmus lépéseinek vezérlése}
    \label{fig:algorithm-controls}
\end{figure}

A játéktábla alatti állapotsor a következőket jeleníti meg:
\begin{compactitem}
    \item Az algoritmus által végrehajtott aktuális \textbf{lépésszám} (step count).
    \item A \textbf{megoldó állapota} (Inicializálva, Fut, Megoldva, Nincs megoldás, Zsákutcába jutott).
    \item A \textbf{megoldás hossza} (lépések száma), ha talált az algoritmus megoldást.
    \item Az éppen megjelenített állapot \textbf{heurisztikus értéke}.
\end{compactitem}

Amikor az algoritmus megoldást talál, a megoldási útvonal zölddel kiemelve jelenik meg az állapotgráfon, és láthatóvá válik a megoldás hossza.


\section{Állapotok felfedezése}

A keresőalgoritmusokon kívül az alkalmazás \textbf{BFS} (szélességi keresés) és \textbf{DFS} (mélységi keresés) felfedező módokat is kínál, amelyek a Solver settings fülről érhetők el. Ezek a módok anélkül tárják fel az állapotteret, hogy megoldást keresnének, minden elérhető állapotot és élt hozzáadva az állapotgráfhoz:

\begin{compactitem}
    \item Beállítható egy \textbf{maximális állapotszám} a felfedezett állapotok korlátozására. Ez ajánlott a nehezebb feladványoknál, ahol az állapottér több ezer elemet is tartalmazhat, mivel a túl sok csúcs átláthatatlanná teheti az állapotgráfot, lelassíthatja a fizikai szimulációt.
    \item Az \textbf{unlimited} kapcsoló eltávolítja az állapotlimitet.
\end{compactitem}

A felfedezés nagy sebességgel fut és feltölti a gráfot, ami után a felhasználó futtathat egy keresőalgoritmust a már teljesen (vagy részben) feltérképezett állapottéren.


\section{Interakció a 3D-s állapotgráffal}

Az alkalmazás bal oldalán található 3D-s térben jeleníti meg az állapotgráfot. Az interakció a következőképpen történik:

\begin{description}
    \item[Kamera forgatása] Tartsa lenyomva a jobb egérgombot és húzza az egeret a gráf körüli keringéshez.
    \item[Kamera mozgatása] Használja a \texttt{W}, \texttt{A}, \texttt{S}, \texttt{D} billentyűket a kamera aktuális irányához képesti eltolásához. \texttt{Shift} nyomva tartásával a mozgás gyorsabb.
    \item[Zoom] Használja az egérgörgőt a közelítéshez és távolításhoz.
    \item[Csúcsra kattintás] Kattintson a bal egérgombbal egy csúcsra az adott állapotba ugráshoz: a játéktábla Kézi módba vált és a kiválasztott állapotot mutatja.
    \item[Csúcsok/élek elrejtése] Egy kapcsológombbal minden csúcs és él szinte teljesen átlátszóvá tehető. Ez hasznos lehet a vizuális zaj csökkentésére nagy gráfok felfedezésekor.
\end{description}


\section{Színmagyarázat}
\label{sec:color-legend}

A 3D-s gráf csúcsai és élei színkódoltak a szerepük szemléltetésére:

\begin{table}[H]
    \centering
    \begin{tabular}{ | m{0.20\textwidth} | m{0.15\textwidth} | m{0.55\textwidth} | }
        \hline
        \textbf{Elem} & \textbf{Szín} & \textbf{Jelentés} \\
        \hline \hline
        \emph{Csúcs} & Piros & Manuálisan kiválasztott aktuális állapot \\
        \hline
        \emph{Csúcs} & Zöld & Megoldott állapot (főautó a kijáratnál) \\
        \hline
        \emph{Csúcs} & Kék & Az aktuális feladvány kezdőállapota \\
        \hline
        \emph{Csúcs} & Narancs & Az algoritmus aktuális állapota \\
        \hline
        \emph{Csúcs} & Sárga & Az algoritmus memóriájában tárolt állapot \\
        \hline
        \emph{Csúcs} & Fehér & Felfedezett állapot \\
        \hline
        \emph{Él} & Zöld & A megtalált megoldási útvonal része \\
        \hline
        \emph{Él} & Sárga (áttetsző) & Az algoritmus memóriájában tárolt lépés \\
        \hline
        \emph{Él} & Fehér (halvány) & Felfedezett lépés \\
        \hline
    \end{tabular}
    \caption{A 3D-s állapotgráf csúcsainak és éleinek színmagyarázata}
    \label{tab:color-legend}
\end{table}

\begin{figure}[H]
    \centering
    \begin{subfigure}[t]{0.48\textwidth}
        \centering
        \includegraphics[width=\linewidth]{yellow_states}
        \caption{Gráfkereső algoritmus memóriájában tárolt feszítőfa látszik sárga színnel, az aktuális állapot narancssárgán van ábrázolva.}
    \end{subfigure}\hfill
    \begin{subfigure}[t]{0.48\textwidth}
        \centering
        \includegraphics[width=\linewidth]{green_states}
        \caption{Megoldott állapot zölddel kitüntetve, hozzá vezető zöld élek mentén a megoldás, a megoldási úton egy állapot van kiválasztva piros színnel megjelenítve.}
    \end{subfigure}
    \caption{Állapotgráf színkódolás szemléltetése}
    \label{fig:color-states}
\end{figure}

Egy állapot csúcsára egyszerre több tulajdonság is vonatkozhat, például épp kiválasztott és algoritmus memóriájában tárolt állapot. Ekkor a táblázatban felsorolt sorrend szerint a legelső tulajdonság megjelenítése élvez prioritást.
Az első négy tulajdonság (kézi aktuális, megoldott, kezdő, algoritmus aktuális) nem csak színben mások, de kétszeres méretben jelennek meg. Az algoritmus memóriájában tárolt csúcsok másfélszeres méretben vannak megjelenítve.
